\documentclass[11pt]{article}
\usepackage[a4paper,margin=28mm]{geometry}
\usepackage[T1]{fontenc}
\usepackage{lmodern,microtype,amsmath,amssymb,amsthm,booktabs,tabularx,array}
\usepackage[colorlinks=true,allcolors=blue!60!black]{hyperref}
\usepackage{xcolor}
\hypersetup{pdftitle={The 2n-1 arc minimum in Koljancic--Oblak, Problem 1.8},pdfauthor={Haobo Ma; Wenlin Zhang}}
\newtheorem{theorem}{Theorem}
\newtheorem{lemma}[theorem]{Lemma}
\newtheorem{corollary}[theorem]{Corollary}
\newcommand{\R}{\mathbb R}
\newcommand{\Zc}{\mathcal Z}
\newcommand{\nssp}{\mathrm{nSSP}}
\newcommand{\leanlink}[1]{\href{https://github.com/the-omega-institute/trureturing/blob/cc684e6326490f9173bfefcadb0294ee6ccbb0b5/D5/S3/Combinatorics/SpectralGraph/KoljancicOblakMinimumArcs.lean}{\texttt{#1}}}
\setlength{\emergencystretch}{2em}
\title{The $2n-1$ arc minimum in Koljan\v{c}i\'{c}--Oblak,\\Problem 1.8 (p.~759)}
\author{Haobo Ma\\\small ChronoAI Pte Ltd, Singapore
\and Wenlin Zhang\\\small National University of Singapore\\\small ChronoAI Pte Ltd, Singapore}
\date{November 2, 2026}
\begin{document}
\maketitle
\begin{abstract}
Koljan\v{c}i\'{c}--Oblak's Problem 1.8, \emph{Electronic Journal of Linear
Algebra} \textbf{42} (2026), p.~759, has an affirmative answer: the minimum
number of arcs in an $n$-vertex digraph requiring the nonsymmetric strong
spectral property is $2n-1$ for every $n\geq1$. By their identification
with Arav et al.'s Problem 6.2, this also gives the minimum $2n-1$ of
nonzero entries in a zero--nonzero pattern requiring the
similarity-transversality property (STP).
We give a self-contained exposition of the centralizer--incidence proof
formalized in Lean \cite{Formal}. More generally, every real matrix with this property and
with $c$ weak support components has at least $2n-c$ nonzero entries.
For connected supports, even allowing one such matrix forces $2n-1$
arcs. Together with the endpoint-looped double path of
Koljan\v{c}i\'{c}--Oblak, this also answers both questions in their
Problem 2.9 (p.~764). The endpoint construction is prior work.
\end{abstract}

\section{The question and its quantifiers}
Koljan\v{c}i\'{c}--Oblak \cite[Problem 1.8, p.~759]{KO} ask for the least
number of arcs in an $n$-vertex digraph \emph{requiring} $\nssp$.
We prove that the answer is $2n-1$, including $n=1$.
Their p.~759 identifies this question as an equivalent reformulation of
Arav et al.'s Problem 6.2 \cite{Arav}. Their p.~756 cites the same work
for the equivalence of $\nssp$ and the similarity-transversality property
(STP). Thus the theorem also answers that zero--nonzero pattern question;
the problem number and equivalence here follow the identification in
\cite[pp.~756, 759]{KO}.
Here arcs are ordered pairs in $[n]\times[n]$, loops count once,
and parallel arcs are absent. For $E\subseteq[n]\times[n]$, put
\[
 \mathcal M(E)=\{A\in\R^{n\times n}:a_{ij}\ne0\ \Longleftrightarrow\ (i,j)\in E\}.
\]
A real matrix $A$ has $\nssp$ if
\begin{equation}\label{eq:nssp}
 A\circ X=0,\qquad AX^{\mathsf T}=X^{\mathsf T}A
 \quad\Longrightarrow\quad X=0
 \qquad(X\in\R^{n\times n}).
\end{equation}
There is no additional diagonal or trace condition. A support
\emph{requires} $\nssp$ if every $A\in\mathcal M(E)$ has it; it
\emph{allows} $\nssp$ if at least one does. These are the conventions
of \cite[Definitions 1.1 and 1.3]{KO}. Signs and nonzero arc weights
are unrestricted throughout.

\begin{theorem}\label{thm:main}
For every $n\geq1$, the minimum of $|E|$ over supports requiring
$\nssp$ is $2n-1$. The same minimum is $2n-1$ when $E$ is restricted
to strongly connected supports, whether it is required to require
$\nssp$ or merely to allow it.
\end{theorem}

The universal requiring result is the theorem encoded by
\leanlink{result} in \cite{Formal}. The argument below also makes its
connected-allowing consequence explicit. This exposition combines that
centralizer proof with the source's endpoint construction.

\section{A centralizer inside the directed cycle space}
Let $\Zc(A)=\{F:AF=FA\}$ be the real centralizer. Write $m=|E|$ and
let $B\in\R^{n\times m}$ be the signed incidence matrix whose column
for $(i,j)$ is $e_j-e_i$. A loop therefore contributes a zero column.

\begin{lemma}\label{lem:centralizer}
Every real $n\times n$ matrix satisfies $\dim_\R\Zc(A)\geq n$.
\end{lemma}
\begin{proof}
Give $V=\R^n$ its $\R[t]$-module structure $t\cdot v=Av$.
The cyclic decomposition theorem for finitely generated torsion modules
over the principal ideal domain $\R[t]$ writes
$V\simeq\bigoplus_r\R[t]/(p_r)$, with $\sum_r\deg p_r=n$.
On each summand, multiplication by an arbitrary residue class commutes
with multiplication by $t$. Extending it by zero on the other summands
gives commuting real linear maps. The resulting map
$\bigoplus_r\R[t]/(p_r)\longrightarrow\Zc(A)$ is injective:
on the $r$th summand its value at the residue class of $1$ recovers the
multiplier. Its domain has real dimension $n$.
\end{proof}

\begin{theorem}\label{thm:cycle}
If $A\in\mathcal M(E)$ has $\nssp$ and $E$ has $c$ weak components,
including isolated vertices, then
\begin{equation}\label{eq:bound}
 n\ \leq\ \dim_\R\Zc(A)\ \leq\ \dim\ker B
   =m-n+c,\qquad m\geq2n-c.
\end{equation}
\end{theorem}
\begin{proof}
Define the linear map
\[
 T_A:\Zc(A)\longrightarrow\R^E,\qquad
       (T_AF)_{(i,j)}=a_{ij}f_{ji}.
\]
If $T_AF=0$, then $A\circ F^{\mathsf T}=0$, including the entries
outside $E$. Since $AF=FA$, apply \eqref{eq:nssp} to $X=F^{\mathsf T}$
to obtain $F=0$. Thus $T_A$ is injective.
For each vertex $v$,
\[
 (BT_AF)_v=\sum_i a_{iv}f_{vi}-\sum_j a_{vj}f_{jv}
         =(FA)_{vv}-(AF)_{vv}=0.
\]
Hence $T_A(\Zc(A))\subseteq\ker B$. Finally,
$\ker B^{\mathsf T}$ consists exactly of functions constant on each
weak component: its equation on arc $(i,j)$ is $p_j-p_i=0$.
It has dimension $c$, so $\operatorname{rank}B=n-c$ and
$\dim\ker B=m-n+c$. Lemma~\ref{lem:centralizer} finishes the proof.
\end{proof}

\begin{corollary}\label{cor:connected}
A weakly connected support allowing $\nssp$ has at least $2n-1$
arcs. Consequently no sign pattern on a weakly connected support with
at most $2n-2$ nonzero positions allows $\nssp$.
If a connected $\nssp$ matrix has exactly $2n-1$ nonzero entries,
then $\dim\Zc(A)=n$ and $T_A:\Zc(A)\to\ker B$ is an isomorphism.
\end{corollary}
\begin{proof}
Set $c=1$ in \eqref{eq:bound}. A matrix in any sign class is still a
matrix with that exact support. In the equality case both inequalities
in \eqref{eq:bound} are equalities; the injective map $T_A$ has domain
and codomain of the same dimension.
\end{proof}

\section{Why universal requiring supports are connected}
We supply the connectivity step rather than assuming it.

\begin{lemma}\label{lem:realize}
Every finite real exact support containing a loop has a realization
with eigenvalue $1$.
\end{lemma}
\begin{proof}
Choose a loop $(r,r)$ and let $N$ have entry $1$ at every support
position except $(r,r)$, and entry $0$ elsewhere. Set
$A(\epsilon,t)=\epsilon N+t e_re_r^{\mathsf T}$.
The determinant is affine in $t$:
\[
 \det(I-A(\epsilon,t))=u(\epsilon)-t v(\epsilon),
 \qquad u(0)=v(0)=1.
\]
Here $u(\epsilon)=\det(I-\epsilon N)$ and $v(\epsilon)$ is the
$(r,r)$ cofactor of $I-\epsilon N$, including the empty cofactor $1$
for a one-vertex support. These are real polynomials. For sufficiently
small nonzero real $\epsilon$, both remain nonzero. Taking
$t=u(\epsilon)/v(\epsilon)$ gives precisely the required support and
$\det(I-A)=0$.
\end{proof}

\begin{lemma}\label{lem:connected}
A support on a nonempty vertex set requiring $\nssp$ is weakly connected.
\end{lemma}
\begin{proof}
Each weak component must contain a loop. Otherwise its diagonal
indicator matrix $F$ is nonzero, commutes with every realization $A$,
and satisfies $A\circ F^{\mathsf T}=0$, contradicting $\nssp$.
If there are two or more components, partition them into two nonempty
groups. Each group contains a loop. By Lemma~\ref{lem:realize}, choose
exact-support matrices $C,D$ on the two groups, both with eigenvalue
$1$. Choose nonzero real vectors $u,v$ with $Cu=u$ and
$v^{\mathsf T}D=v^{\mathsf T}$. Then, for
\[
 A=\begin{pmatrix}C&0\\0&D\end{pmatrix},\qquad
 F=\begin{pmatrix}0&uv^{\mathsf T}\\0&0\end{pmatrix},
\]
one has $AF=FA$ and $F\ne0$. Since there are no arcs between the
groups, $A\circ F^{\mathsf T}=0$. This contradicts the requirement
for this realization.
\end{proof}

The distinction from allowing is essential: $\operatorname{diag}(1,2)$
has $\nssp$ and two arcs, while its disconnected support does not
require $\nssp$, as the realization $I_2$ shows. In particular the
connected-support qualification in Corollary~\ref{cor:connected}
cannot be removed.

\section{The sharp example, with its arbitrary-weight proof}
Let $P_{n,\{1\}}$ have both arcs between consecutive vertices and the
single loop $(1,1)$. It has $2(n-1)+1=2n-1$ arcs and is strongly
connected, with the usual length-zero-path convention at $n=1$.
The fact that this support requires $\nssp$ is
\cite[Theorem 5.5 with $m=1$, p.~771; Remark 2.10(A), p.~764]{KO}.
Here is a direct proof, using the odd-offset propagation also present
in \cite{Formal}.

\begin{lemma}\label{lem:path}
Every real matrix with exact support $P_{n,\{1\}}$ has $\nssp$.
\end{lemma}
\begin{proof}
The case $n=1$ is immediate. Otherwise write
$a_{11}=d\ne0$, $a_{i,i+1}=u_i\ne0$ and
$a_{i+1,i}=\ell_i\ne0$ for $1\leq i<n$; all other entries vanish.
Let $AF=FA$ and $A\circ F^{\mathsf T}=0$. Then
$f_{11}=0$ and $f_{i,i+1}=f_{i+1,i}=0$.
Use zero entries outside the index interval $[1,n]$, and put
$\ell_n=0$.

We first prove that every entry at odd distance from the diagonal is
zero. The first off-diagonals already vanish. Suppose $k\geq3$ is odd
and all smaller odd offsets vanish. In the $(i+1,i+k)$ entry of
$AF=FA$, both diagonal coefficients of $A$ are zero. The terms at
offset $k-2$ vanish, leaving
\begin{equation}\label{eq:odd}
 \ell_i f_{i,i+k}=\ell_{i+k}f_{i+1,i+k+1}
 \qquad(1\leq i\leq n-k).
\end{equation}
Starting with $i=n-k$, the right side is zero. Descend in $i$ and
use $\ell_i\ne0$ to obtain the entire upper offset $k$. Applying the
same argument to $(A^{\mathsf T},F^{\mathsf T})$ proves the lower
offset. Induction establishes the claim.

It follows that the first column of $F$ is zero. Its first entry is
already zero, and all even-indexed entries have odd offset. For odd
$i\geq3$, the $(i,1)$ equation of $AF=FA$ reduces to
$0=d f_{i1}$: the neighboring entries in column $1$ and $f_{i2}$
have odd offset, and $a_{ii}=0$. Thus these entries also vanish.
Finally, the vectors $e_1,Ae_1,\ldots,A^{n-1}e_1$ form a basis:
$A^ke_1$ is supported on the first $k+1$ coordinates and its
$(k+1)$st coordinate is $\ell_1\cdots\ell_k\ne0$.
Since $F$ commutes with $A$ and $Fe_1=0$, it kills this basis, so
$F=0$. Setting $F=X^{\mathsf T}$ proves \eqref{eq:nssp}.
\end{proof}

\begin{proof}[Proof of Theorem~\ref{thm:main}]
By Lemma~\ref{lem:connected}, a requiring support is weakly connected.
Corollary~\ref{cor:connected} gives its lower bound $2n-1$.
Lemma~\ref{lem:path} supplies equality at the same order $n$.
For the two strongly connected variants, the lower bound follows
already from connected allowing, while the same path provides
attainment and even the stronger requiring property.
\end{proof}

\section{Source coverage and verification scope}
The following table distinguishes consequences of this proof from
earlier results and questions outside its scope. All locators in the
first column refer to the journal version \cite{KO} unless specified.
\begin{center}\small
\begin{tabularx}{\textwidth}{@{}p{0.30\textwidth}X@{}}
\toprule
Source item & Coverage\\
\midrule
Problem 1.8, p.~759 & Proved: the requiring minimum is $2n-1$.\\
Problem 2.9, p.~764, both questions & Proved: both strongly connected minima
(requiring and allowing) are $2n-1$.\\
Theorem 2.7(B), p.~763 & Generalized: any weakly connected support with at
most $2n-2$ arcs fails to allow $\nssp$, not just supports over trees.\\
Remark 2.8, p.~764 & Its finite-order requiring conclusions follow for every
positive order.\\
Theorem 5.5, p.~771; Remark 2.10(A), p.~764 & Prior construction credited;
the endpoint-loop case is proved above. Other initial loop segments are
not rederived.\\
Problem 1.6, p.~757 & Full require/allow/not-allow trichotomy not covered.
Prior allow/not-allow classification: \cite[Theorem 6.1]{Paths}.\\
Question 5.11, p.~776 & Prior affirmative proof:
\cite[Corollary 6.2]{Paths}, for even $n$ and every nonempty $L$.\\
Arav et al.\ \cite{Arav}, Problem 6.2, as identified in
\cite[p.~759]{KO} & Proved: the minimum number of nonzero entries in a
zero--nonzero pattern requiring STP is $2n-1$; the equivalence with $\nssp$
is recorded in \cite[p.~756]{KO}.\\
Arav et al.\ \cite{Arav}, Problems 6.1 and 6.5, as identified in
\cite[pp.~757--759]{KO} & Prior negative answers: Example 1.7 disproves the
unrestricted sign minimum in 6.1; p.~757 disproves 6.5.\\
Arav et al.\ \cite{Arav}, Problem 6.3, as identified in
\cite[p.~763]{KO} & General sign-pattern allowability classification
not covered here; \cite[Theorem 2.7]{KO} gives particular obstructions.\\
\bottomrule
\end{tabularx}
\end{center}
The connected sign-pattern consequence is Corollary~\ref{cor:connected},
not a repair of the unrestricted sign question: its four-entry
counterexample in \cite[Example 1.7]{KO} is disconnected.
The allow/not-allow classification in \cite{Paths} is a dated public
mathematical note, not a journal publication. Its statements are
credited here rather than presented as new consequences of our bound.
The classification of all equality supports is not established here.

\paragraph{Formal correspondence.}
Every link below points directly to the same pinned Lean file
at commit \texttt{cc684e6326490f9173bfefcadb0294ee6ccbb0b5}.
All declaration names have namespace
\texttt{D5.S3.Combinatorics.SpectralGraph.}\allowbreak
\texttt{KoljancicOblakMinimumArcs}.
\begin{center}\small
\begin{tabularx}{\textwidth}{@{}p{0.55\textwidth}X@{}}
\toprule
Lean declaration & Mathematical statement\\
\midrule
\leanlink{nssp\_centralizer\_cycle\_bound} &
$n+\operatorname{rank}B\leq |E|$ for any real exact-support $\nssp$ matrix.\\
\leanlink{nssp\_arc\_bound\_of\_connected} &
$|E|\geq2n-1$ for a connected support allowing $\nssp$.\\
\leanlink{requires\_connected} & A requiring support is weakly connected.\\
\leanlink{endpoint\_requires}, \leanlink{endpoint\_attainment} &
The endpoint-looped double path requires $\nssp$ and has $2n-1$ arcs.\\
\leanlink{result} & The universal all-positive-order requiring minimum.\\
\bottomrule
\end{tabularx}
\end{center}
The main theorem and its mathematical deductions have written proofs above. The $c$-component
form, the equality isomorphism and the explicit source-coverage
corollaries are written deductions; they are not additional named
declarations in that file. The accompanying \texttt{verify.py} checks
finite integer matrices by exact rational and finite-field arithmetic.
Its finite ranges are stated in \texttt{verification.json}; the
universal statements follow from the proofs, not from those samples.

\begin{thebibliography}{9}
\bibitem{KO}
S. Koljan\v{c}i\'{c} and P. Oblak,
\emph{Combinatorial aspects of the nonsymmetric strong spectral property
for graphs}, Electronic Journal of Linear Algebra \textbf{42} (2026),
755--779. \href{https://doi.org/10.13001/ela.2026.10523}{doi:10.13001/ela.2026.10523}.
Versioned preprint: \href{https://arxiv.org/abs/2604.21552v1}{arXiv:2604.21552v1}.
\bibitem{Paths}
Carptopus, \emph{Spectral separation and the non-symmetric strong spectral
property for looped double paths}, MathExplore Notes, v0.1-beta,
19 August 2026, Theorem 6.1 and Corollary 6.2.
\href{https://github.com/Carptopus/MathExplore-Notes/blob/e57ee007380cbc351eb371d8e991303cab5f1e48/research/nssp-looped-double-paths/nssp-looped-double-paths.tex}{Commit-pinned mathematical text}.
\bibitem{Formal}
\emph{KoljancicOblakMinimumArcs}, in \texttt{trureturing},
commit \texttt{cc684e6326490f9173bfefcadb0294ee6ccbb0b5},
\leanlink{Lean source}.
\bibitem{Arav}
M. Arav, F.\,J. Hall, H. van der Holst, Z. Li, A. Mathivanan,
J. Pan, H. Xu and Z. Yang,
\emph{Advances on similarity via transversal intersection of manifolds},
Linear Algebra and its Applications \textbf{721} (2025), 102--121.
\href{https://doi.org/10.1016/j.laa.2024.05.013}{doi:10.1016/j.laa.2024.05.013}.
\end{thebibliography}
\medskip
\noindent\small\textcopyright\ 2026 Haobo Ma and Wenlin Zhang.
This note is licensed under
\href{https://creativecommons.org/licenses/by/4.0/}{Creative Commons Attribution 4.0 International (CC BY 4.0)}.
\end{document}
